\section{Characterization techniques}
\label{sec:techniques}
Each card lists: \textbf{principle} (how it works), \textbf{what you can do with it}, \textbf{limits}, and the \textbf{exercises}
that need it.

\newcommand{\tech}[5]{%
\begin{cardbox}{#1}
\textbf{Principle:} #2\par\smallskip
\textbf{You can:} #3\par
\textbf{Limits:} #4\par
\textbf{Exercises:} #5
\end{cardbox}}

\subsection{UV-vis and fluorescence}
\tech{UV-vis \slide{124--136}}
{UV and visible photons (\SIrange{200}{800}{nm}) have the energy of electronic transitions, so the sample absorbs at wavelengths
set by its electronic structure. In molecules and semiconductors these are HOMO$\to$LUMO or band-gap transitions. In Au, Ag and Cu
nanoparticles the light's electric field drives the conduction electrons to oscillate collectively against the ionic cores: the
\textbf{localised surface plasmon resonance (LSPR)}. At resonance absorption is very strong and gives a band in the visible. Its
position depends on the metal, size, shape, composition and surroundings of the particle. The colour we see is the complementary
colour of the absorbed light (Au absorbs green and looks red; Ag absorbs violet and looks yellow). Larger particles have electrons
less bound to the atoms, so the oscillation is ``softer'' and the band moves to longer $\lambda$; in rods the electrons can also
oscillate along the long axis, which gives a second, lower-energy band.}
{\begin{itemize}
 \item identify the \textbf{metal}: Ag $\approx$ 400 nm, Au $\approx$ 520 nm, Cu $\approx$ 570 nm; oxides show no band;
 \item tell an \textbf{alloy} (one band in between, moving with composition) from segregated metals (two bands);
 \item see \textbf{size} (bigger $\to$ red shift + broader; small effect below $\sim$30 nm) and \textbf{distribution} (band width);
 \item see \textbf{shape}: rods give a transverse ($\sim$520 nm) and a longitudinal band that red-shifts with aspect ratio.
\end{itemize}}
{indirect; only plasmonic metals/semiconductors; tiny particles ($\lesssim$3 nm) show a damped band. Fast, cheap and
representative of the whole sample.}
{1, 10d, 13, 14, 15b, 16.}
\textbf{Fluorescence/quantum dots:} the sample absorbs a photon, relaxes non-radiatively within the excited state and then emits a
photon of lower energy (Stokes shift). In CdSe quantum dots the band gap grows as the size decreases (quantum confinement), so small
dots emit blue and large ones red \slide{133--135}. Used for \emph{in vivo} imaging.

\subsection{X-ray photoelectron spectroscopy (XPS)}
\tech{XPS \slide{137--165}}
{X-ray photons (Al K$\alpha$ 1486.6 eV, Mg K$\alpha$ 1253.6 eV) eject core electrons by the photoelectric effect. A hemispherical
analyser measures their kinetic energy, and energy conservation gives the binding energy:
$E_\mathrm{kin} = h\nu - E_B - \phi$ ($\phi$ = spectrometer work function). $E_B$ is characteristic of each element and orbital.
It also shifts by a few eV with the chemical environment: core electrons feel the nucleus screened by all other electrons, so
removing valence charge (oxidation) reduces the screening and raises $E_B$ (the \textbf{chemical shift}). X-rays penetrate microns,
but photoelectrons lose energy after travelling only one inelastic mean free path (0.5--3 nm), so only those from the top
\textbf{5--10 nm} leave without losses. This makes XPS a \textbf{surface technique}. It needs ultra-high vacuum so the electrons
reach the detector.}
{\begin{itemize}
 \item identify elements at the surface; quantify with $I_i = N_i\sigma_i\lambda_iK$ (peak areas, Shirley background);
 \item determine \textbf{oxidation state} from the chemical shift: losing valence charge (oxidation, more electronegative
       neighbour) $\Rightarrow E_B\uparrow$; gaining charge $\Rightarrow E_B\downarrow$ (Ti$^0$ 454.1 $\to$ TiO$_2$ 458.8 eV);
 \item detect \textbf{electron transfer between two metals}: the more electronegative one gains charge ($E_B\downarrow$), the other
       loses it ($E_B\uparrow$);
 \item compare surface vs core: a shell hides the core; particles outside a tube give more signal than particles inside.
\end{itemize}}
{surface only, UHV (except near-ambient-pressure XPS at synchrotrons, where tunable $h\nu$ also allows depth profiling); small
shifts for Cu$^0$/Cu$^+$ and Ag$^0$/Ag$^+$.}
{4, 5, 9b, 10a, 15c--d, 16.}
\begin{keybox}[: reading peaks]
$s$ levels: one peak. $p$, $d$, $f$: \textbf{doublets} from spin--orbit coupling ($j = l\pm\frac12$) with area ratios
$p_{1/2}:p_{3/2}=1:2$, $d_{3/2}:d_{5/2}=2:3$, $f_{5/2}:f_{7/2}=3:4$; the higher-$j$ peak is at lower $E_B$. Auger peaks (e.g.\ Cu KLL)
have a kinetic energy independent of $h\nu$. Cu$^{2+}$/Ni$^{2+}$ show shake-up satellites\extra. The $E_B$ axis runs from high (left)
to low (right).
\end{keybox}

\subsection{X-ray diffraction (XRD)}
\tech{XRD \slide{167--187}}
{A crystal is a stack of parallel atomic planes, labelled by Miller indices $(hkl)$ with spacing $d_{hkl}$. X-rays (Cu K$\alpha$,
1.54 \AA) reflected by successive planes interfere constructively only when their path difference is a whole number of wavelengths:
\textbf{Bragg's law} $2d\sin\theta = n\lambda$. Each family of planes therefore gives a peak at its own angle $\twotheta$; for cubic
crystals $d_{hkl} = a/\sqrt{h^2+k^2+l^2}$. In a powder, randomly oriented crystallites produce all allowed reflections at once.
In centred lattices some reflections cancel because an extra plane of atoms sits halfway and scatters in antiphase (systematic
absences), so the list of observed $(hkl)$ reveals the lattice type. A finite crystal cannot cancel the intensity fully slightly away
from the Bragg angle, so \textbf{small crystallites give broad peaks} (Scherrer).}
{\begin{itemize}
 \item \textbf{identify phases} (compare expected and measured $\twotheta$);
 \item get the \textbf{lattice parameter} and the \textbf{lattice type} (fcc/bcc) by indexing;
 \item distinguish \textbf{mixture} (two sets of peaks) from \textbf{solid solution} (one set at intermediate positions, Vegard\extra);
 \item estimate \textbf{crystallite size} with Scherrer, $L = 0.9\lambda/(B\cos\theta)$, $B$ = FWHM $-$ instrumental, in radians.
\end{itemize}}
{only crystalline phases; very small/amorphous particles give broad or no peaks; minor phases ($<$ few \%) invisible; cannot tell
Au from Ag (lattice parameters 0.3\,\% apart). Crystallite size $\neq$ particle size (a particle can contain several
crystallites).}
{2, 3, 4, 5, 6, 16.}
\begin{keybox}[: selection rules \slide{179--182}]
SC: all $hkl$. \textbf{BCC}: $h+k+l$ even $\to$ 110, 200, 211, 220, 310, 222 ($N$ = 2, 4, 6, 8, 10, 12; ratios 1:2:3:4:5:6).
\textbf{FCC}: all odd or all even $\to$ 111, 200, 220, 311, 222, 400 ($N$ = 3, 4, 8, 11, 12, 16; ratios 1:1.33:2.67:3.67:4:5.33,
peaks in pairs). Rocksalt (TiN, NiO) and fluorite (CeO$_2$) follow the fcc rule.
\end{keybox}

\subsection{Transmission electron microscopy (TEM)}
\tech{TEM, HRTEM, STEM, EDX, EELS \slide{188--218}}
{A beam of high-energy electrons passes through a very thin sample (particles dispersed on a carbon-coated Cu grid). Electrons
scattered elastically (same energy, new direction) form images and diffraction patterns; inelastically scattered ones (energy lost)
carry chemical information. An aperture in the back focal plane of the objective lens selects the electrons used for the image.
\textbf{Bright field} uses the direct beam, so heavy or thick regions that scatter strongly look dark. \textbf{Dark field/HAADF}
collects scattered electrons, whose intensity grows as $Z^2$, so heavy atoms look bright (Z-contrast). In \textbf{HRTEM} the direct
and diffracted beams interfere and give \textbf{lattice fringes} spaced by $d_{hkl}$ (an interference pattern, not images of atoms);
their Fourier transform (FFT) looks like a diffraction pattern. The beam also excites characteristic X-rays (\textbf{EDX}) and
element-specific energy losses (\textbf{EELS}); filtering the image at one energy loss gives elemental maps (\textbf{EFTEM}).}
{\begin{itemize}
 \item measure \textbf{size distribution}: $\bar d = \sum n_id_i/\sum n_i$ with $\geq$100 particles; see shape and location;
 \item identify a particle or phase from fringe spacings, and the \textbf{exposed facets} (fringes parallel to the surface);
 \item judge \textbf{epitaxy}: fringes of particle and support aligned across the interface;
 \item map \textbf{which element is where} (core/shell, alloy) with EDX/EELS; see heavy metals on light supports with HAADF;
 \item find the 3D location (inside/outside a tube or pore) with \textbf{tomography}; follow reactions in situ (environmental TEM).
\end{itemize}}
{local (few particles, not representative; one big particle can hold 50\,\% of the mass \slide{202}), 2D projection, vacuum, beam
damage, organics invisible.}
{5, 7, 8, 9, 10b--c, 13, 14, 15b--c, 16, 17, 18.}

\subsection{Scanning probe microscopy (AFM, STM)}
\tech{AFM and STM \slide{219--237}}
{A sharp tip is scanned over the surface by a piezoelectric scanner with sub-nm precision, and a feedback loop keeps the tip--sample
interaction constant. The image is the record of the tip's height. \textbf{AFM}: the tip sits on a flexible cantilever; the force
from the surface bends it, and a laser reflected onto a photodiode measures the bending (contact mode, or tapping mode with an
oscillating cantilever for soft samples). The result is the real \textbf{topography} of anything on the surface. \textbf{STM}: a
metal tip a few \AA{} from a conducting sample lets electrons tunnel through the vacuum gap, $I_t\sim V\rho(E_F)e^{-2kd}$ with
$k\approx\SI{1}{\per\angstrom}$, so 1 \AA{} more distance lowers $I_t$ about tenfold. This extreme sensitivity gives atomic
resolution. In constant-current mode the tip follows the local density of electronic states, not just the topography.}
{\begin{itemize}
 \item AFM: measure precise \textbf{heights} (sub-nm) on any material, in air or liquid, \textbf{including organic shells};
 \item STM: get \textbf{atomic resolution} and the local density of states on conducting surfaces; move single atoms (IBM Xe logo,
       Fe quantum corrals).
\end{itemize}}
{AFM lateral sizes overestimated by the tip shape; STM needs conducting, clean, flat samples (UHV).}
{17 (capped particle on a surface: $h_\mathrm{AFM}\approx d_\mathrm{TEM}+2\ell_\mathrm{cap}$); 15c (SEM/AFM see only outside).}

\subsection{Which technique for which question}
\begin{center}
\small
\renewcommand{\arraystretch}{1.15}
\setlength{\tabcolsep}{3pt}
\begin{tabularx}{\linewidth}{@{}>{\raggedright\arraybackslash}p{4.4cm}X>{\raggedright\arraybackslash}p{2.2cm}@{}}
\toprule
Question & Best techniques (what you see) & Exercises \\ \midrule
Particle size & TEM histogram; XRD Scherrer (crystallite); AFM height; UV-vis (indirect) & 4, 6, 7, 14, 17 \\
Phase, lattice parameter & XRD; HRTEM fringes (local) & 2--6, 9, 10, 16 \\
Mixture vs solid solution & XRD (two sets vs one shifted set); STEM-EDX maps; HRTEM per particle & 6, 16 \\
Metal vs oxide & XPS (shift, satellites, O 1s); XRD (structure); UV-vis (plasmon only for metal) & 4, 5, 10 \\
Core/shell, alloy & EELS/EDX maps; UV-vis (one vs two bands); XPS before/after shell removal; HAADF & 9, 13, 16 \\
Electronic interaction & XPS shifts + electronegativity & 15d, 16 \\
Exposed facets, epitaxy & HRTEM + FFT & 5, 8, 10c, 15b, 18 \\
Inside vs outside & TEM tomography; XPS intensity; SEM/AFM & 15c \\
Capping shell present & AFM height vs TEM core; XPS (C 1s, S 2p) & 17 \\
Whole-sample average & UV-vis, XRD (not TEM) & --- \\
\bottomrule
\end{tabularx}
\end{center}
